\section{Introduction}

\subsection{Urban context and motivation}
\begin{bluenote}
Introduce the problem of traffic disruptions (planned works, incidents, etc.) and explain why
selecting detours is difficult due to network topology and traffic conditions. It motivates the need
for a decision-support approach that can compare mutliple recovery plans in a realistic way rather
than relying on intuitions and simple metrics.
\end{bluenote}

\subsubsection{Digital Twin}

A digital twin is a data-driven virtual representation of real-world entities designed to mirror a
physical object or system throughout its lifecycle. The virtual part is updated with real-time data
gotten from sensors or other sources making it possible to link the physical and digital parts.
\cite{defDT,IBMDT}

A digital twin consists of three components \cite{DT_health}:
\begin{enumerate}
\item \textbf{Physical entity}: The real-world object or system being modeled (e.g., traffic, a
plant)
\item \textbf{Virtual representation}: A model that replicates the behavior of the physical entity.
\item \textbf{Connection}: The link between the physical system and the virtual system making it
possible to update the physical system by exchanging information via this real-time connection.
\end{enumerate}

\begin{figure}[htbp]
\centering
\safeincludegraphics[width=0.7\linewidth]{chapters/introduction/pictures/digital_twin.png}
\caption{Illustration of a digital twin from \cite{digital_twin}}
\label{enter-label}
\end{figure}

Digital twins receive continuously streams of data from sensors or other sources ensuring that the
virtual model replicates the current state of the physical object. The main goals for a digital twin
are to simulate the physical entity and predict what could happen in order to optimize or predict
failures. Urban traffic networks, which are both data-rich and fast-evolving, provide fertile ground
for the deployment of digital twins (see \thesisref{dt_traffic_twin}). Digital twins in traffic can
potentially simulate and optimize traffic flows, supporting smarter city planning and real-time
operations. A detailed discussion of modern digital twin solutions in traffic management is provided
in \thesisref{dt_traffic_simulation}.

\begin{figure}[htbp]
\centering
\safeincludegraphics[width=0.8\linewidth]{chapters/introduction/pictures/sumo_twin.png}
\caption{xample workflow for driving SUMO towards a digital twin. Various static and real-time data
sources (maps, sensors, cameras) are integrated into simulation models to continuously update and
calibrate the digital twin for real-world application. Source : \cite{sumo_twin}}
\label{dt_traffic_twin}
\end{figure}

\subsubsection{Context}

In this project, a digital twin acts as a precise and efficient method to simulate traffic based on
real-world data. Digital twins leverage continuous data streams (for example, from vehicle counting
devices) to create an up-to-date virtual model of physical traffic flows. Accurate traffic
simulation relies on high-resolution data, typically gathered by a variety of counting devices such
as inductive loop traffic detectors, Bluetooth recognition systems and ANPR (Automatic Number Plate
Recognition) cameras \cite{counting_devices}. These devices measure the number of vehicles passing
specific locations over time, producing the data needed for dynamic modeling and simulation.

Through the aggregation and analysis of such data, it becomes possible to estimate traffic density
and infer vehicle movement patterns across a road network. In this framework, physical traffic
corresponds to the physical entity of the digital twin, real-time vehicle counts represent the
virtual layer and the counting devices themselves serve as the communication link between the
physical and digital components.

Traffic simulation is particularly vital in urban environments such as Brussels, where the frequent
congestion and ongoing roadworks demand dynamic management of the transport network. Planned
deviations due to construction may alter typical traffic flows and produce unexpected bottlenecks.
Accurate simulation enables city planners to evaluate alternative deviation plans and identify those
that minimize congestion and disruption. Additionally, the direct and indirect effects of traffic
interventions, such as road closures, can be promptly assessed.

\subsection{Problem statement}
\begin{bluenote}
The thesis problem : Given a baseline traffic situation and a disruptive event (road closure) the
goal is to generate feasible deviation plans and assess their impact on traffic in order to
recommend the best plan to the urban planner.
\end{bluenote}

\subsubsection{Objective}

The main objective of this work is the development of a web-based application that reproduces and
extends the functionalities of an existing tool, TrafficTwin\thesisurlfootnote{https://traffictwin.ulb.be}
\cite{SIM_param}. This modernization enhances accessibility and usability. The work is divided into
several sub-goals:

\begin{rednote}
TO BE WRITTEN
\end{rednote}

Ultimately, the goal is to provide city planners and traffic managers with actionable, data-driven
tools for congestion mitigation and urban mobility planning.

\subsection{Research questions}
\begin{bluenote}
Present the research questions that will be addressed in the thesis. The main one is ``How
deviation plans should be evaluated?'' and the subquestions are ``What are the key factors
influencing the evaluation of deviation plans?'' and ``How can we develop a decision-support
approach to compare multiple deviation plans?''.
\end{bluenote}

\subsection{Contributions and thesis outline}
\begin{bluenote}
Summarize the contributions (end-to-end deviation plan assessment pipeline, candidate plan
generation strategies, Monte Carlo evaluation protocol). It ends with a short outline of the thesis
chapters.
\end{bluenote}
